\chapter[Sermon 79. A Doleful Song or Mourning, and Lamentation of Rome, Which the Princes and Merchant Make for Her.]{A Doleful Song or Mourning, and Lamentation of Rome, Which the Princes and Merchant Make for Her.}
\label{sermon:079}
\markright{Sermon 79. A Doleful Song or Mourning, and Lamentation of Rome ...}

And the kings of the earth shall weep and wail over her, those who committed adultery with her and lived sensually, when they see the smoke of her burning, and shall stand afar off, for fear of her punishment, saying, “Alas, alas, that great city Babylon, that mighty city! For in one hour your judgment has come.” And the merchants of the earth shall weep and wail in themselves, because no one will buy their wares anymore—the wares of gold and silver, and precious stones, of pearls, and silk, and purple, and scarlet, and all manner of wood, and all manner of vessels of ivory, and all manner of vessels of most precious wood, and of brass, and of iron, and cinnamon, and perfumes, and ointments, and frankincense, and wine, and oil, and fine flour, and livestock, and sheep, and horses, and chariots, and bodies and souls of men. And the delights that your soul desired have departed from you. And all things which were luxurious and were of great price have departed from you, and you shall find them no more. The merchants of these things, who became rich from her, shall stand afar off for fear of her punishment, weeping and wailing, saying, “Alas, alas, that great city, that was clothed in silk and purple, and scarlet, and decked with gold and precious stones and pearls! For in one hour so great riches have come to nothing.” And every ship’s captain, and all who occupy ships, and the sailors who work on the sea, stood afar off and cried when they saw the smoke of her burning, and said, “What city is like this great city?” And they cast dust on their heads, and cried weeping and wailing, and said, “Alas, alas, the great city, in which all who had ships on the sea were made rich by reason of her wares! For in one hour she is made desolate.”

\index{Jerome, Saint}\index{Ezekiel (Book of)}\index{Rome}Fourthly, this chapter addresses the weeping, lamentation, or mourning of Rome, which will be struck down and destroyed. The text is abundant and marvelous, employing a clear hypostasis that presents all things before our eyes. Our Lord God consistently uses a familiar method: whenever he intends to clearly demonstrate and impress upon all people the downfall or destruction of a nation, kingdom, or city, he commands his prophets to compose an elegy, or a lamentable song. Such lamentations reveal not only the subversion but also the causes of destruction, and the manner of desolation is recounted. The purpose is also declared, so that others may not become like that nation and share in its destruction. We have clear examples in the writings of the prophets, especially the lamentations of Jerome, and the mournful song of Tyre in Ezekiel, chapters 27 and 28, which aligns more closely with this passage. Indeed, it appears that John has borrowed many things from that source.

There is nothing here to greatly concern ourselves with. The sum of it all is this: Rome will fall and perish completely, so that nothing remains, neither of the empire nor of that See, much less of the riches and pleasures. This was partly fulfilled in ancient Rome, and partly will be fulfilled in the new Rome at the day of judgment.

However, neither Christ himself, nor the Apostle is recorded lamenting the downfall of Babylon, but rather wicked persons are induced [?]. These are Kings and Princes of the Earth, merchants, or governors of ships, or mariners, who have all committed fornication with this harlot and, through her company, have become rich. Indeed, old Rome was supplied with the friendship of Kings; and again, the officials sent by them to govern provinces appeared to be Kings and Princes. And because the riches of Rome were great, and all nations were plunged into riot, the merchants there gained exceedingly. Moreover, ships sailed to Rome from the East, South, and West—that is, from Syria, Egypt or Africa, and from Spain itself, and the furthest parts of the world. But when Rome was destroyed and lay in ruins, and the Empire was torn apart, those whose profit and pleasure were lost could do nothing but lament.

\index{Antichrist}\index{Erasmus of Rotterdam}New Rome has also, besides those temporal, even peculiar merchants and princes of her own. For the prelates of the church are princes. And in the church of Rome, all the saints of God know, how to occupy the traffic of merchandise. For what holy thing is not to be bought in that seat? Merchandise is practiced in forgiveness of sins, in pardons and satisfactions, in ecclesiastical benefices, in the worship of images and saints, in masses, in burials, in saying diriges for the dead, and almost in all spiritual matters. Hereof comes an immeasurable gain, and the greatest occasion of pleasures. Other merchants buy their wares very dearly; the Roman Canaanites pay not one denier or farthing for their wares, but sell them for an unreasonable price. Neither suppose I that ever there was any merchandise like unto this in all the world, nor yet a more gainful profit from a thing of naught. Erasmus has also touched these things, in the proverb to ask tribute of a dead man. And whereas before the day of judgment, the Lord Christ shall destroy Antichrist with the spirit of his mouth, and that gain begins to be diminished, we see how everywhere among these spiritual merchants, complaints and grudgings arise. Then what manner of lamentation and wailing think you will be, where the same Lord by his coming shall utterly abolish the same Antichrist, and they must go into fire everlasting?

Again, we must somewhat also consider the mourning. To mourn of itself, is no sin. For the best and holiest men have lamented their dead, and their calamities, and destruction of cities and realms. For Abraham mourned. The lamentations also of Jeremiah remain, over the city of Jerusalem. The faithful mourned with a great mourning for Stephen in the Acts. However, in their lamentation they kept a mean, and referred all things to the glory of God, and salvation of souls. The ungodly and worldly men do not mourn after this sort. They never remember the signs of men, for which the righteous Lord punishes the world, neither do they refer the evils of them and theirs to the glory, verity, and justice of God, or amendment of manners: therefore are they not sorry that God is offended, nor require forgiveness of sins; but it grieves them that occasion of sinning is taken from them, that their pleasures and profit is past.

And now, while princes, merchants, and mariners—not for the favor of God lost, not of true compassion or love of their neighbor, but for love of themselves, for the loss of earthly things, for the destruction of goodly, ancient, strong, and precious things—but chiefly for their profit lost and pleasures taken away. The Apostle makes mention of such grief in the second to the Corinthians, chapter 7. And surely this sorrow and mourning is nothing else but a description and a shadowing of a most certain and greatest destruction, and that of men ungodly. And full well and purposefully doth he set forth the wailing both in the behavior of the mourners, and also by their words. To their gesture appertains that they weep, wail, cry out, and cast dust on their heads. To their words are referred these things: “Woe, woe, alas, alas, that great city!” \&c., which is repeated of the merchants and sailors.

Moreover, here are also touched upon the causes of destruction, the riot and luxury in which Rome indulged. And likewise are rehearsed the wealth, riches, majesty, pride, and pleasures of Rome. And here, by the way, are warned all worldly men regarding what they may expect if they devote themselves to the pleasures and luxury of this world, which was present in Rome and is immeasurable. We have not read in any histories that nations have long endured, whether they have been blessed or conquered by worldly pleasures. To build, to eat, to drink, to be clothed, and to have servants, men and women, is lawful; but a measure must be kept in these things, as in all others. The benefits of God must be acknowledged, and they may not be valued more highly than virtue. But at Rome, and in the world, godliness and moderation are passed over, and these things alone are regarded, desired, and beloved. In buildings and household goods, all things were sumptuous and immeasurable. They were made of gold, which could have been well of earth or tin; of silver, where wood or iron might have served. And when wood was chosen, it was not every wood, but the most excellent was chosen. [illegible] appears to be named of Thyia, a tree, to which Theophrastus attributes great honor, reporting that the famous buildings of old temples were made thereof, and a certain immortality of matter, incorrupt, enduring on houses against all weathers, etc. Pliny has this in Book 13, Chapter 16. In service also, they use men like beasts; neither do they have beasts for their own use, but mostly chosen ones. They have horses and mules exceeding fine. They have their horse litters, coaches, and chariots, truly notable; all things glisten with gold, precious stones, and purple; and all things are wrought and devised for pride and sumptuousness. What shall we say that the whole bands of their men go all in silks and velvet, wearing their masters’ colors? The lord himself of all, sitting on the shoulders of his Belphougers, is borne on high and is carried on men’s bodies as the most noble chariot. In the meat and drink of these men, all things are most delicate, exquisite, and varied. Their drink is costly, strange, and immoderate. The apparel of their body is also over-sumptuous. Their garments glisten with gold and are stiff with pearls. Their common garment is of crimson satin. They use also ointments and apples of desire, which may both be understood of the fruits of trees, and also of pomanders containing musk and smelling sweet, and of odoriferous savors.

Finally, in all things, we must consider what the ultimate fate is of riot, pride, and sensuality, and how unstable the favor and friendship of men are. Here, everything perishes; nothing remains secure. And those who have prepared for many years perish truly in one hour. They flee from us in danger, having received great gain from our hands. Indeed, they stand far off, out of danger, and lament the sorrowful circumstance; no one comes to help or deliver us. Every man is afraid of his own life. Let us therefore learn to trust in God, to despise pleasures, to place no confidence in flesh or the friendship of men. For while you are fortunate, you will have many friends; if the world begins to frown upon you, they will all forsake you, in whom you placed your trust, and leave you in the thorns. And this is the chief end of all these things, as I showed at the beginning: Rome shall fall and be made desolate forever. The Lord our God restrain all evil. Amen.
